\input{../tex-inputs/latex-leseansicht-vorspann}

\section[Carl von Torresani an Arthur Schnitzler, 6. 2. 1892]{L04327 Carl von Torresani an Arthur Schnitzler, 6. 2. 1892}
\nopagebreak\mylabel{L04327v}
\rehead{ }\normalsize\beginnumbering\briefempfaengerindex{Schnitzler, Arthur@\textsc{Schnitzler, Arthur}!zzzTorresani-Lanzenfeld, Carl von@\emph{von Carl von Torresani-Lanzenfeld}!1892-02-061@{6. 2. 1892}|(be}
\toendnotes[C]{\smallbreak\pagebreak[2]}
\correspDesc{Versand  durch Carl von Torresani am 6. 2. 1892 in Paris
\newline{}Erhalt  durch Arthur Schnitzler im Zeitraum [7. 2. 1892
                  – 11. 2. 1892?] in Wien}\toendnotes[C]{\smallbreak}
\Standort{CUL, Schnitzler, B 106.}
\physDesc{Brief, 2 Blätter, 6 Seiten, 3843 Zeichen
\newline{}Handschrift: schwarze Tinte, deutsche Kurrent
\newline{}Schnitzler: mit Bleistift beschriftet: »\textsc{Torresani}« }\toendnotes[C]{\smallbreak}
\pstart
           \raggedleft{}{\pb}Paris\oindex{Paris@\textbf{Paris}, \emph{Hauptstadt}|pw} am 6. Februar 1892.\pend
           
\pstart{}Hochgeſchätzter Herrn Doktor!\pend\vspace{0.5em}
\pstart
           Wenn ich für \label{K_L04327-1v}\edtext{Ihre liebe Sendung\pwindex{Schnitzler, Arthur 15.\,5.\,1862 Wien – 21.\,10.\,1931 ebd.@\textsc{Schnitzler, Arthur} (15.\,5.\,1862 Wien – 21.\,10.\,1931 ebd.), \emph{Schriftsteller, Mediziner}!Märchen. Schauspiel in drei Aufzügen@\strich\emph{Das Märchen. Schauspiel in drei Aufzügen}|pwv}}{\lemma{\textnormal{\emph{Ihre liebe Sendung}}}\Cendnote{\textnormal{Schnitzlers Sendung, die ein
                  Widmungsexemplar von \emph{Das Märchen}\pwindex{Schnitzler, Arthur 15.\,5.\,1862 Wien – 21.\,10.\,1931 ebd.@\textsc{Schnitzler, Arthur} (15.\,5.\,1862 Wien – 21.\,10.\,1931 ebd.), \emph{Schriftsteller, Mediziner}!Märchen. Schauspiel in drei Aufzügen@\strich\emph{Das Märchen. Schauspiel in drei Aufzügen}|pwk} enthielt,
                  ist nicht überliefert. Schnitzler bekam die
                  Drucke des Stücks\pwindex{Schnitzler, Arthur 15.\,5.\,1862 Wien – 21.\,10.\,1931 ebd.@\textsc{Schnitzler, Arthur} (15.\,5.\,1862 Wien – 21.\,10.\,1931 ebd.), \emph{Schriftsteller, Mediziner}!Märchen. Schauspiel in drei Aufzügen@\strich\emph{Das Märchen. Schauspiel in drei Aufzügen}|pwkv} am 5. 12. 1891 und sandte
                  noch am selben Tag eines an Hugo August von
                     Hofmannsthal\pwindex{Hofmannsthal, Hugo August von 21.\,12.\,1841 Wien – 8.\,12.\,1915 ebd.@\textsc{Hofmannsthal, Hugo August von} (21.\,12.\,1841 Wien – 8.\,12.\,1915 ebd.), \emph{Bankdirektor}|pwk}, siehe XXXX Auszeichnungsfehler: Dokument L00048 nicht gefunden.}}}\label{K_L04327-1} erſt heute danke, ſo bitte ich es nicht meiner Faulheit, noch weniger
               meiner Indifferenz zuſchreiben zu wollen. Ich kann sie ehrlich verſichern, daß \introOben{}ſie\introOben{} mir eine herzliche Freude u Überraſchung bereitet hat.
               Auch habe ich das Stück\pwindex{Schnitzler, Arthur 15.\,5.\,1862 Wien – 21.\,10.\,1931 ebd.@\textsc{Schnitzler, Arthur} (15.\,5.\,1862 Wien – 21.\,10.\,1931 ebd.), \emph{Schriftsteller, Mediziner}!Märchen. Schauspiel in drei Aufzügen@\strich\emph{Das Märchen. Schauspiel in drei Aufzügen}|pwv}{ }ſofort
               mit größtem Intereſſe geleſen, acht Tage ſpäter nochmals überleſen; und nur mit dem
               Brief bin ich in \introOben{}einer\introOben{} Verſpätung, die sie meiner jetzt ganz
               beſonders ſtarken Arbeitsüberhäufung zu Gute halten werden.\pend
           
\pstart
           Nehmen Sie, verehrter Herr Doktor, meine warme Gratulation; Sie haben da eine
               prächtige Leiſtung zuwege ge{\pb}bracht. Ich
               will nicht lobhudeln; \strikeout{A} es würde Sie von meiner Seite
               gewiß nicht freuen, denn ich glaube, Sie \strikeout{hab} erwarten
               von mir nicht fade Complimente. Aber ich kann nur ſagen: dieſe fein\strikeout{e} gedachte, unendlich zart durchgeführte,
               pſychologiſch tiefe u jeder Effekthaſcherei aus dem Wege gehende Comödie\pwindex{Schnitzler, Arthur 15.\,5.\,1862 Wien – 21.\,10.\,1931 ebd.@\textsc{Schnitzler, Arthur} (15.\,5.\,1862 Wien – 21.\,10.\,1931 ebd.), \emph{Schriftsteller, Mediziner}!Märchen. Schauspiel in drei Aufzügen@\strich\emph{Das Märchen. Schauspiel in drei Aufzügen}|pwv} muß \strikeout{ſ\textcolor{gray}{×}\-\textcolor{gray}{×}\-\textcolor{gray}{×}} unter die beſten gerechnet werden, welche in dieſer Art exiſtiren. Der
               Charakter der Fanny\pwindex{Schnitzler, Arthur 15.\,5.\,1862 Wien – 21.\,10.\,1931 ebd.@\textsc{Schnitzler, Arthur} (15.\,5.\,1862 Wien – 21.\,10.\,1931 ebd.), \emph{Schriftsteller, Mediziner}!Märchen. Schauspiel in drei Aufzügen@\strich\emph{Das Märchen. Schauspiel in drei Aufzügen}|pwv} iſt
               geradezu meiſterhaft gezeichnet, und dabei durch und durch wahr. Jener des Fedor\pwindex{Schnitzler, Arthur 15.\,5.\,1862 Wien – 21.\,10.\,1931 ebd.@\textsc{Schnitzler, Arthur} (15.\,5.\,1862 Wien – 21.\,10.\,1931 ebd.), \emph{Schriftsteller, Mediziner}!Märchen. Schauspiel in drei Aufzügen@\strich\emph{Das Märchen. Schauspiel in drei Aufzügen}|pwv} hat für mich einige
               Unverſtändlichkeiten, die ich mir{ }ſeinerzeit von Ihnen, der ihn gewiſs in jeder Weiſe
               überdacht hat, mündlich erklären laſſen werde; im ganzen u Großen erſcheint er mir
               als ein moderner Hamlet\pwindex{\textcolor{red}{\textsuperscript{XXXX indx1}}!Hamlet@\strich\emph{Hamlet}|pwv}, voll
                  Fein{\pb}heit u voll Intereſſe. Die Mutter\pwindex{Schnitzler, Arthur 15.\,5.\,1862 Wien – 21.\,10.\,1931 ebd.@\textsc{Schnitzler, Arthur} (15.\,5.\,1862 Wien – 21.\,10.\,1931 ebd.), \emph{Schriftsteller, Mediziner}!Märchen. Schauspiel in drei Aufzügen@\strich\emph{Das Märchen. Schauspiel in drei Aufzügen}|pwv}{ }ſehe ich leibhaftig vor
               mir; ich kenne ſie, ſo zu ſagen. Die Nebencharaktere der beiden Gigerln{ }ſind voll
               humaniſtiſcher, discret gezeichneter Pointen und von einer Unaufdringlichkeit, die
               ich bewundern muß. \uline{Ich} hätte mich abſolut zu{ }ſtärkerem Auftragen der Farben verleiten laſſen; Sie haben die Charakterſtärke
               gehabt, dieſe Klippe zu umſchiffen. Aus Wandel\pwindex{Schnitzler, Arthur 15.\,5.\,1862 Wien – 21.\,10.\,1931 ebd.@\textsc{Schnitzler, Arthur} (15.\,5.\,1862 Wien – 21.\,10.\,1931 ebd.), \emph{Schriftsteller, Mediziner}!Märchen. Schauspiel in drei Aufzügen@\strich\emph{Das Märchen. Schauspiel in drei Aufzügen}|pwv} werde ich nicht recht klug; der Mann iſt offenbar
               ein \introOben{}kleinherziger\introOben{}{ }\label{K_L04327-2v}\edtext{Biedermeier}{\lemma{\textnormal{\emph{Biedermeier}}}\Cendnote{\textnormal{hier abwertend für ›Biedermann‹ (mit einer Anspielung auf
                   die Epoche des Biedermeiers), also
                   konservativ, spießbürgerlich und kleingeistig.}}}\label{K_L04327-2}; dabei hat
               er aber gewiſſe edle Züge, welche Lächerlichkeit ausſchließt. Ich glaube Sie haben
               ihn eben nach der Natur gezeichnet, objectiv, ohne Voreingeno{\geminationm}enheit. Kurz, ich kann nur mit Überzeugung
               applaudiren.\pend
           
\pstart
           Jetzt erlauben Sie mir aber zwei unverzeihliche kritiſche Bemerkungen, oder beſſer
               Fragen. 1.) Glauben Sie nicht, {\pb}daß
               etwas zu viele \uline{junge Leute} auftreten? Werden{ }ſie
               nicht, gleich alt, gleich (ſchwarz) gekleidet, alſo nothwendigerweiſe mit einer
               gewiſſen gegenſeitigen Ähnlichkeit behaftet, im Publicum eine Art Confuſion
               hervorbringen, das{ }ſich Zurechtfinden erſchweren? Ich glaube, Sie hätten mindeſtens
               zwei \strikeout{\textcolor{gray}{×}\-\textcolor{gray}{×}\-\textcolor{gray}{×}\-\textcolor{gray}{×}\-\textcolor{gray}{×}\-\textcolor{gray}{×}\-\textcolor{gray}{×} können}, insbeſonders den
               Bruder des D\textsuperscript{r}{ }\textsc{Witte\pwindex{Schnitzler, Arthur 15.\,5.\,1862 Wien – 21.\,10.\,1931 ebd.@\textsc{Schnitzler, Arthur} (15.\,5.\,1862 Wien – 21.\,10.\,1931 ebd.), \emph{Schriftsteller, Mediziner}!Märchen. Schauspiel in drei Aufzügen@\strich\emph{Das Märchen. Schauspiel in drei Aufzügen}|pwv}} eliminieren u \strikeout{Ihr} ihre Rolle in die anderer
               verweben könne{[}n{]}.\pend
           
\pstart
           2.) Sind Sie überzeugt, daß das \introOben{}(große)\introOben{} Theaterpublicum die
               Feinheit des Stückes\pwindex{Schnitzler, Arthur 15.\,5.\,1862 Wien – 21.\,10.\,1931 ebd.@\textsc{Schnitzler, Arthur} (15.\,5.\,1862 Wien – 21.\,10.\,1931 ebd.), \emph{Schriftsteller, Mediziner}!Märchen. Schauspiel in drei Aufzügen@\strich\emph{Das Märchen. Schauspiel in drei Aufzügen}|pwv}, die
               vornehme \uline{Enthaltſamkeit} des Autors auffaſſen kann?
               Ich kenne nicht das Publicum des \label{K_L04327-3v}\edtext{Leſſingtheaters\oindex{Lessing-Theater@\textbf{Lessing-Theater}, \emph{Theater}|pw}}{\lemma{\textnormal{\emph{Lessingtheaters}}}\Cendnote{\textnormal{Schnitzler hatte brieflich am XXXX Auszeichnungsfehler: Dokument L00054 nicht gefunden durch Richard Beer-Hofmann\pwindex{Beer-Hofmann, Richard 11.\,7.\,1866 Wien – 26.\,9.\,1945 New York City@\textsc{Beer-Hofmann, Richard} (11.\,7.\,1866 Wien – 26.\,9.\,1945 New York City), \emph{Schriftsteller}|pwk} erfahren, dass das Stück\pwindex{Schnitzler, Arthur 15.\,5.\,1862 Wien – 21.\,10.\,1931 ebd.@\textsc{Schnitzler, Arthur} (15.\,5.\,1862 Wien – 21.\,10.\,1931 ebd.), \emph{Schriftsteller, Mediziner}!Märchen. Schauspiel in drei Aufzügen@\strich\emph{Das Märchen. Schauspiel in drei Aufzügen}|pwkv} vom Berliner\oindex{Berlin@\textbf{Berlin}, \emph{Hauptstadt}|pwk}{ }\emph{Lessingtheater}\orgindex{Lessing-Theater@Lessing-Theater|pwk} angenommen sei. Die
                  Inszenierung verzögerte sich jedoch und unterblieb bis 1925. Die Uraufführung\eventindex{Volkstheater@\textbf{Volkstheater}!Uraufführung von Das Märchen, 1.12.1893@Uraufführung von Das Märchen, 1.12.1893|pwkv} fand stattdessen am
                   1. 12. 1893 am \emph{Volkstheater}\orgindex{Volkstheater@Volkstheater|pwk} in
                     Wien\oindex{Wien@\textbf{Wien}, \emph{Verwaltungsgebiet}|pwk} statt.}}}\label{K_L04327-3}, wünſche aber von Herzen,
               es möge ein Elite-Publicum{ }ſein; dann wird der reichlichſte, wärmſte Beifall nicht
               ausbleiben, wie ich ihn Ihnen von Herzen zolle, alſo wünſche.\pend
           
\pstart
           Ihr »Lied des Alkandi\pwindex{Schnitzler, Arthur 15.\,5.\,1862 Wien – 21.\,10.\,1931 ebd.@\textsc{Schnitzler, Arthur} (15.\,5.\,1862 Wien – 21.\,10.\,1931 ebd.), \emph{Schriftsteller, Mediziner}!Alkandi’s Lied@\strich\emph{Alkandi’s Lied}|pw}« iſt mir \uline{ans Herz gewachſen}, doch {\pb}darüber mündlich nächſten Winter, in
               welchem ich mir vornehme, Ihre{ }ſympathiſche, leider all zu{ }ſpät gemachte
               Bekanntſchaft recht gründlich zu pflegen.\pend
           
\pstart
           F. Salten\pwindex{Salten, Felix 6.\,9.\,1869 Budapest – 8.\,10.\,1945 Zürich@\textsc{Salten, Felix} (6.\,9.\,1869 Budapest – 8.\,10.\,1945 Zürich), \emph{Schriftsteller, Journalist, Chefredakteur}|pw} hat mir wiederholt über Sie
               geſchrieben u damit einem regen Wunſche meinerseits Rechnung getragen. Er iſt auch
               ein charmanter u wertvoller Menſch. Könnten Sie mir nicht{ }ſeine jetzige Adreſſe{ }ſchicken? Es iſt aus der \label{K_L04327-4v}\edtext{Redaction\orgindex{Allgemeine Kunst-Chronik@Allgemeine Kunst-Chronik|pwv}}{\lemma{\textnormal{\emph{Redaction}}}\Cendnote{\textnormal{Felix
                     Salten\pwindex{Salten, Felix 6.\,9.\,1869 Budapest – 8.\,10.\,1945 Zürich@\textsc{Salten, Felix} (6.\,9.\,1869 Budapest – 8.\,10.\,1945 Zürich), \emph{Schriftsteller, Journalist, Chefredakteur}|pwk} arbeitete 1890 und 1891 für die \emph{Allgemeine Kunst-Chronik}\orgindex{Allgemeine Kunst-Chronik@Allgemeine Kunst-Chronik|pwk} und besprach dort
                  unter Kürzel vier Bücher\pwindex{Torresani-Lanzenfeld, Carl von 19.\,4.\,1846 Mailand – 16.\,4.\,1907 Nago-Torbole@\textsc{Torresani-Lanzenfeld, Carl von} (19.\,4.\,1846 Mailand – 16.\,4.\,1907 Nago-Torbole), \emph{Schriftsteller, Offizier}!Mit tausend Masten. Roman@\strich\emph{Mit tausend Masten. Roman}|pwkv}\pwindex{Torresani-Lanzenfeld, Carl von 19.\,4.\,1846 Mailand – 16.\,4.\,1907 Nago-Torbole@\textsc{Torresani-Lanzenfeld, Carl von} (19.\,4.\,1846 Mailand – 16.\,4.\,1907 Nago-Torbole), \emph{Schriftsteller, Offizier}!Auf gerettetem Kahn. Roman@\strich\emph{Auf gerettetem Kahn. Roman}|pwkv}\pwindex{Salten, Felix 6.\,9.\,1869 Budapest – 8.\,10.\,1945 Zürich@\textsc{Salten, Felix} (6.\,9.\,1869 Budapest – 8.\,10.\,1945 Zürich), \emph{Schriftsteller, Journalist, Chefredakteur}!Juckerkomtesse«@\strich\emph{»Die Juckerkomtesse«}|pwkv}\pwindex{Torresani-Lanzenfeld, Carl von 19.\,4.\,1846 Mailand – 16.\,4.\,1907 Nago-Torbole@\textsc{Torresani-Lanzenfeld, Carl von} (19.\,4.\,1846 Mailand – 16.\,4.\,1907 Nago-Torbole), \emph{Schriftsteller, Offizier}!Schwarzgelbe Reitergeschichten@\strich\emph{Schwarzgelbe Reitergeschichten}|pwkv} von Torresani\pwindex{Torresani-Lanzenfeld, Carl von 19.\,4.\,1846 Mailand – 16.\,4.\,1907 Nago-Torbole@\textsc{Torresani-Lanzenfeld, Carl von} (19.\,4.\,1846 Mailand – 16.\,4.\,1907 Nago-Torbole), \emph{Schriftsteller, Offizier}|pwk}: \emph{»Schwarzgelbe Reitergeschichten«}\pwindex{Salten, Felix 6.\,9.\,1869 Budapest – 8.\,10.\,1945 Zürich@\textsc{Salten, Felix} (6.\,9.\,1869 Budapest – 8.\,10.\,1945 Zürich), \emph{Schriftsteller, Journalist, Chefredakteur}!Schwarzgelbe Reitergeschichten«@\strich\emph{»Schwarzgelbe Reitergeschichten«}|pwk}. In: \emph{Allgemeine Kunst-Chronik}\pwindex{Allgemeine Kunst-Chronik@\emph{Allgemeine Kunst-Chronik}|pwk}, XIV. Jg., H. 26
                     (zweites Dezemberheft 1890), S. 735–736; \emph{»Mit tausend Masken.« – »Auf gerettetem
                        Kahn«}\pwindex{Salten, Felix 6.\,9.\,1869 Budapest – 8.\,10.\,1945 Zürich@\textsc{Salten, Felix} (6.\,9.\,1869 Budapest – 8.\,10.\,1945 Zürich), \emph{Schriftsteller, Journalist, Chefredakteur}!Mit tausend Masken.« – »Auf gerettetem Kahn.«@\strich\emph{»Mit tausend Masken.« – »Auf gerettetem Kahn.«}|pwk}. In: ebd.\pwindex{Allgemeine Kunst-Chronik@\emph{Allgemeine Kunst-Chronik}|pwkv},
                     XV. Jg., H. 4 (erstes Februarheft 1891),
                     S. 112–114 und \emph{»Die Juckerkomtesse«}\pwindex{Salten, Felix 6.\,9.\,1869 Budapest – 8.\,10.\,1945 Zürich@\textsc{Salten, Felix} (6.\,9.\,1869 Budapest – 8.\,10.\,1945 Zürich), \emph{Schriftsteller, Journalist, Chefredakteur}!Juckerkomtesse«@\strich\emph{»Die Juckerkomtesse«}|pwk}. In: ebd.\pwindex{Allgemeine Kunst-Chronik@\emph{Allgemeine Kunst-Chronik}|pwkv}, H. 26 (zweites
                        Dezemberheft 1891), S. 720–721.}}}\label{K_L04327-4}
               ausgetreten u hat vergeſſen mir mitzutheilen, wie ich ihm{ }ſchreiben{ }ſoll.\pend
           
\pstart
           Ich arbeite noch immer an einem Roman\pwindex{Torresani-Lanzenfeld, Carl von 19.\,4.\,1846 Mailand – 16.\,4.\,1907 Nago-Torbole@\textsc{Torresani-Lanzenfeld, Carl von} (19.\,4.\,1846 Mailand – 16.\,4.\,1907 Nago-Torbole), \emph{Schriftsteller, Offizier}!Kinder der Weltstadt. Wiener Roman@\strich\emph{Kinder der Weltstadt. Wiener Roman}|pwv}, der \label{K_L04327-5v}\edtext{im
                  März}{\lemma{\textnormal{\emph{im
                  März}}}\Cendnote{\textnormal{Ab dem Ostersonntag 1892
                  erschien in 66 Folgen Karl Baron Torresani\pwindex{Torresani-Lanzenfeld, Carl von 19.\,4.\,1846 Mailand – 16.\,4.\,1907 Nago-Torbole@\textsc{Torresani-Lanzenfeld, Carl von} (19.\,4.\,1846 Mailand – 16.\,4.\,1907 Nago-Torbole), \emph{Schriftsteller, Offizier}|pwk}: \emph{Kinder der Weltstadt. Wiener Roman}\pwindex{Torresani-Lanzenfeld, Carl von 19.\,4.\,1846 Mailand – 16.\,4.\,1907 Nago-Torbole@\textsc{Torresani-Lanzenfeld, Carl von} (19.\,4.\,1846 Mailand – 16.\,4.\,1907 Nago-Torbole), \emph{Schriftsteller, Offizier}!Kinder der Weltstadt. Wiener Roman@\strich\emph{Kinder der Weltstadt. Wiener Roman}|pwk}. In: \emph{Neues Wiener Tagblatt}\pwindex{Neues Wiener Tagblatt@\emph{Neues Wiener Tagblatt}|pwk},
                        17. 4. 1892, Nr. 108 bis 24. 6. 1892,
                     Nr. 174.}}}\label{K_L04327-5} im neuen Wr Tagblatt\pwindex{Neues Wiener Tagblatt@\emph{Neues Wiener Tagblatt}|pw} erſcheint. Er wird
               mir{ }ſchwer, denn ich betrete damit ein ganz neues Feld, die ewige Ariſtokraterei iſt
               mir für den Augenblick langweilig geworden, u. ich will nicht warten, bis ſie es auch
               dem Publicum wird.\pend
           
\pstart
           {\pb}Nehmen Sie nochmals, beſter Herr Doktor,
               meinen herzlichen Dank für Buch\pwindex{Schnitzler, Arthur 15.\,5.\,1862 Wien – 21.\,10.\,1931 ebd.@\textsc{Schnitzler, Arthur} (15.\,5.\,1862 Wien – 21.\,10.\,1931 ebd.), \emph{Schriftsteller, Mediziner}!Märchen. Schauspiel in drei Aufzügen@\strich\emph{Das Märchen. Schauspiel in drei Aufzügen}|pwv} u Widmung, u die Verſicherung \uline{wahrſter}
               Verehrung entgegen, womit ich bin\pend
           
\pstart
           Ihr ganz ergebner{\\[\baselineskip]}\spacefill\mbox{B. Torresani}\pend
           \leftskip=0em{}
\pstart
           \noindent{}\textsc{Paris, 18, R. Clément Marot\oindex{rue Clément-Marot, 18@\textbf{rue Clément-Marot, 18}, \emph{Wohngebäude}|pw}}.\pend
           \selectlanguage{ngerman}\endnumbering\briefempfaengerindex{Schnitzler, Arthur@\textsc{Schnitzler, Arthur}!zzzTorresani-Lanzenfeld, Carl von@\emph{von Carl von Torresani-Lanzenfeld}!1892-02-061@{6. 2. 1892}|)be}\mylabel{L04327h}
\begin{anhang}
\end{anhang}\newcommand{\dateiname}{L04327}\newcommand{\titel}{Carl von Torresani an Arthur Schnitzler, 6. 2. 1892}\newcommand{\editorInnen}{Herausgegeben von Jahnke, SelmaMüller, Martin Anton}\input{../tex-inputs/latex-leseansicht-abspann}
